\documentclass[12pt, a4paper]{article}

% ── Packages ──────────────────────────────────────────────────────────────────
\usepackage[margin=2.5cm]{geometry}
\usepackage{amsmath, amssymb, amsthm}
\usepackage{graphicx}
\usepackage{booktabs}
\usepackage{hyperref}
\usepackage{xcolor}
\usepackage{listings}
\usepackage{caption}
\usepackage{subcaption}
\usepackage{float}
\usepackage{enumitem}
\usepackage{titlesec}
\usepackage{fancyhdr}
\usepackage{cite}
\usepackage{array}
\usepackage{multirow}
\usepackage{algorithm}
\usepackage{algpseudocode}

% ── Listings style ────────────────────────────────────────────────────────────
\lstset{
  basicstyle=\ttfamily\small,
  keywordstyle=\color{blue!70!black}\bfseries,
  commentstyle=\color{green!50!black}\itshape,
  stringstyle=\color{red!60!black},
  backgroundcolor=\color{gray!8},
  frame=single,
  rulecolor=\color{gray!30},
  breaklines=true,
  numbers=left,
  numberstyle=\tiny\color{gray},
  tabsize=4,
  showstringspaces=false
}

% ── Header / Footer ───────────────────────────────────────────────────────────
\pagestyle{fancy}
\fancyhf{}
\rhead{Lane Detection \& Object Detection AI}
\lhead{\leftmark}
\cfoot{\thepage}

% ── Hyperref setup ────────────────────────────────────────────────────────────
\hypersetup{
  colorlinks=true,
  linkcolor=blue!70!black,
  citecolor=green!50!black,
  urlcolor=cyan!60!black
}

% ── Title ─────────────────────────────────────────────────────────────────────
\title{
  \vspace{-1cm}
  {\Large\textbf{Road Vision AI}} \\[0.4em]
  {\large Lane Detection Using Cascaded U-Net Architecture \\
   with YOLOv8 Object Detection Integration} \\[0.8em]
  \rule{\textwidth}{0.4pt}
}
\author{Akshwin T \\ \small akshwint.2003@gmail.com}
\date{\today}

% ══════════════════════════════════════════════════════════════════════════════
\begin{document}
% ══════════════════════════════════════════════════════════════════════════════

\maketitle
\thispagestyle{empty}
\tableofcontents
\newpage

% ══════════════════════════════════════════════════════════════════════════════
\section{Abstract}
% ══════════════════════════════════════════════════════════════════════════════

This report presents \textit{Road Vision AI}, a deep learning--based road
analysis system combining \textbf{lane detection} via a custom Cascaded
U-Net architecture with \textbf{real-time object detection} via YOLOv8.
The system is deployed as an interactive Streamlit web application supporting
both portrait and landscape road videos, with a batch-processing pipeline
for multi-video workflows.

The lane segmentation model is trained to predict pixel-level binary masks
identifying drivable lane boundaries from road imagery. Post-inference,
predicted masks are rescaled back to the original frame resolution,
eliminating the aspect-ratio distortion inherent in naive 256$\times$256
resizing. The object detection module uses the YOLOv8 nano/small/medium
models to localise vehicles, pedestrians, and traffic signs in the same
frames.

Key contributions include:
\begin{itemize}[noitemsep]
  \item A portrait-video--aware preprocessing pipeline with configurable
        road-region cropping.
  \item Batched inference (8 frames per forward pass) reducing per-frame
        overhead by $\approx$4--5$\times$.
  \item A frame-skip optimisation (1$\times$, 1.5$\times$, 2$\times$,
        3$\times$) for resource-constrained deployment.
  \item A polished, session-state--correct Streamlit UI with persistent
        download buttons.
\end{itemize}

% ══════════════════════════════════════════════════════════════════════════════
\section{Summary of the Selected Lab Option}
% ══════════════════════════════════════════════════════════════════════════════

\subsection{Lab Option: Road Scene Understanding}

The selected project addresses two complementary tasks in autonomous driving
and road safety:

\begin{enumerate}
  \item \textbf{Lane Detection} — semantic segmentation of lane markings
        using a Cascaded Dilated U-Net (CD-UNet).
  \item \textbf{Object Detection} — localisation and classification of
        road actors (vehicles, pedestrians, signs) using YOLOv8.
\end{enumerate}

The rationale for combining both tasks is that safe autonomous navigation
requires simultaneously knowing \emph{where lanes are} and \emph{what
objects are present}. The lane model is a custom-trained network; the
detection model reuses the publicly available YOLOv8 weights (COCO-trained)
to avoid the cost of annotating a detection dataset from scratch.

\subsection{Cascaded Dilated U-Net Architecture}

The backbone is a U-Net variant in which standard convolutions are replaced
by \emph{dilated (atrous) convolutions}. For a 2-D input feature map $X$
with dilation rate $r$, the dilated convolution of kernel $w$ at position
$(i,j)$ is:

\begin{equation}
  (X \ast_r w)(i,j)
  = \sum_{m}\sum_{n} X(i + r \cdot m,\; j + r \cdot n)\; w(m,n)
  \label{eq:dilated_conv}
\end{equation}

Increasing $r$ exponentially enlarges the receptive field without reducing
spatial resolution, which is crucial for detecting thin, elongated lane
markings.

The \emph{cascade} refers to stacking two U-Net sub-networks end-to-end
so that the coarse prediction of the first stage is fed as an additional
input channel to the second stage:

\begin{equation}
  \hat{Y}_2 = f_{\theta_2}\bigl(X,\; \hat{Y}_1\bigr),
  \qquad
  \hat{Y}_1 = f_{\theta_1}(X)
  \label{eq:cascade}
\end{equation}

where $f_{\theta_k}$ is the $k$-th U-Net and $\hat{Y}_k$ is its binary
output mask.

% ══════════════════════════════════════════════════════════════════════════════
\section{Dataset and Setup Description}
% ══════════════════════════════════════════════════════════════════════════════

\subsection{Lane Detection Dataset}

The model was trained on a curated road-scene dataset composed of dashcam
and street-view frames paired with ground-truth binary masks that indicate
lane-pixel regions.

\begin{table}[H]
  \centering
  \caption{Dataset statistics for lane segmentation training}
  \begin{tabular}{lc}
    \toprule
    \textbf{Property} & \textbf{Value} \\
    \midrule
    Input resolution (training) & $256 \times 256$ px \\
    Image channels & RGB (3) \\
    Mask channels & Binary grayscale (1) \\
    Train / Validation split & 80\,\% / 20\,\% \\
    Augmentations & Horizontal flip, brightness jitter \\
    \bottomrule
  \end{tabular}
  \label{tab:dataset}
\end{table}

\subsection{Object Detection Dataset}

YOLOv8 is used with pre-trained COCO weights. The COCO dataset contains:

\begin{itemize}[noitemsep]
  \item 80 object categories (vehicles, persons, animals, everyday objects).
  \item $\approx$118\,000 training images and $\approx$5\,000 validation images.
  \item Bounding-box and instance-segmentation annotations.
\end{itemize}

Relevant COCO classes for road scenes are grouped as shown in
Table~\ref{tab:classes}.

\begin{table}[H]
  \centering
  \caption{COCO class groups used in the Road Scene filter}
  \begin{tabular}{ll}
    \toprule
    \textbf{Group} & \textbf{Classes (IDs)} \\
    \midrule
    Vehicles    & bicycle(1), car(2), motorcycle(3), bus(5), truck(7), boat(8) \\
    Traffic Signs & traffic light(9), stop sign(11), parking meter(12) \\
    People      & person(0) \\
    \bottomrule
  \end{tabular}
  \label{tab:classes}
\end{table}

\subsection{Deployment Environment}

\begin{table}[H]
  \centering
  \caption{Deployment environment specifications}
  \begin{tabular}{ll}
    \toprule
    \textbf{Component} & \textbf{Specification} \\
    \midrule
    Platform           & Streamlit Community Cloud \\
    OS                 & Debian GNU/Linux 13 (Trixie) \\
    Python             & 3.11.15 \\
    TensorFlow         & 2.17.0 (CPU-only) \\
    Keras              & 3.x (bundled with TF 2.17) \\
    ultralytics        & 8.x (YOLOv8) \\
    OpenCV             & 4.x (headless) \\
    \bottomrule
  \end{tabular}
  \label{tab:env}
\end{table}

% ══════════════════════════════════════════════════════════════════════════════
\section{Code Overview}
% ══════════════════════════════════════════════════════════════════════════════

\subsection{System Architecture}

The application is structured into three logical layers:

\begin{enumerate}
  \item \textbf{Model layer} — \texttt{load\_model()} and \texttt{load\_yolo()}
        are decorated with \texttt{@st.cache\_resource} so models are loaded
        once and shared across all sessions.
  \item \textbf{Inference layer} — pure Python functions operating on
        NumPy arrays; no Streamlit calls, fully testable in isolation.
  \item \textbf{UI layer} — Streamlit widgets, session state management,
        and result rendering.
\end{enumerate}

\subsection{Preprocessing Pipeline}

\subsubsection{Portrait-Video Road-Region Crop}

A configurable \texttt{road\_start\_ratio} $\rho \in [0, 0.7]$ discards
the top $\rho$ fraction of each frame before inference:

\begin{equation}
  \text{ROI}(F) = F\bigl[\lfloor \rho H \rfloor : H,\; 0 : W\bigr]
  \label{eq:crop}
\end{equation}

where $H \times W$ is the original frame size. This focuses the model
on the road surface and removes sky, overhead signs, and buildings that
would otherwise consume the model's 256$\times$256 input budget.

\subsubsection{Image Normalisation}

Each ROI is resized to $256 \times 256$ and normalised to $[0, 1]$:

\begin{equation}
  \tilde{X} = \frac{\text{resize}(\text{ROI}(F),\; 256 \times 256)}{255}
  \label{eq:norm}
\end{equation}

\subsubsection{Mask Re-projection}

After inference the predicted 256$\times$256 mask $M_{256}$ is rescaled
back to the ROI dimensions, then embedded into a zero-initialised
full-frame canvas:

\begin{equation}
  M_{\text{full}}[y, x] =
  \begin{cases}
    \text{resize}(M_{256},\; W \times (H - \lfloor\rho H\rfloor))[y - \lfloor\rho H\rfloor, x]
      & \text{if } y \geq \lfloor\rho H\rfloor \\
    0 & \text{otherwise}
  \end{cases}
  \label{eq:reproject}
\end{equation}

\subsection{Loss Function}

The lane segmentation model is trained with \textbf{binary cross-entropy}
(BCE):

\begin{equation}
  \mathcal{L}_{\text{BCE}}
  = -\frac{1}{N} \sum_{i=1}^{N}
    \Bigl[
      y_i \log \hat{p}_i + (1 - y_i) \log(1 - \hat{p}_i)
    \Bigr]
  \label{eq:bce}
\end{equation}

where $y_i \in \{0,1\}$ is the ground-truth label for pixel $i$,
$\hat{p}_i = \sigma(z_i)$ is the predicted probability, and $\sigma$
is the sigmoid activation:

\begin{equation}
  \sigma(z) = \frac{1}{1 + e^{-z}}
  \label{eq:sigmoid}
\end{equation}

\subsection{Evaluation Metrics}

\subsubsection{Intersection over Union (IoU)}

\begin{equation}
  \text{IoU}
  = \frac{|\hat{Y} \cap Y|}{|\hat{Y} \cup Y|}
  = \frac{\text{TP}}{\text{TP} + \text{FP} + \text{FN}}
  \label{eq:iou}
\end{equation}

\subsubsection{Dice Coefficient (F1 for segmentation)}

\begin{equation}
  \text{Dice}
  = \frac{2|\hat{Y} \cap Y|}{|\hat{Y}| + |Y|}
  = \frac{2\,\text{TP}}{2\,\text{TP} + \text{FP} + \text{FN}}
  \label{eq:dice}
\end{equation}

The Dice coefficient is numerically equivalent to the F1-score and is
the primary metric for imbalanced segmentation tasks (lane pixels are
typically $<5\%$ of total pixels).

\subsubsection{Pixel-wise Accuracy}

\begin{equation}
  \text{Accuracy}
  = \frac{\text{TP} + \text{TN}}{\text{TP} + \text{TN} + \text{FP} + \text{FN}}
  \label{eq:acc}
\end{equation}

\subsection{YOLO Detection}

YOLOv8 divides the input into an $S \times S$ grid. Each cell predicts
$B$ bounding boxes; each box is parametrised as $(t_x, t_y, t_w, t_h, t_o)$
where the decoded centre coordinates relative to a grid cell of offset
$(c_x, c_y)$ and anchor size $(p_w, p_h)$ are:

\begin{align}
  b_x &= \sigma(t_x) + c_x \\
  b_y &= \sigma(t_y) + c_y \\
  b_w &= p_w\, e^{t_w}   \\
  b_h &= p_h\, e^{t_h}   \\
  b_o &= \sigma(t_o)
  \label{eq:yolo_decode}
\end{align}

The final confidence score for class $c$ is:

\begin{equation}
  \text{conf}_c = b_o \cdot P(\text{class}_c \mid \text{object})
  \label{eq:yolo_conf}
\end{equation}

Non-Maximum Suppression (NMS) removes overlapping predictions:

\begin{equation}
  \text{keep box } i \iff
  \text{IoU}(B_i, B_j) < \tau_{\text{NMS}}
  \;\;\forall j \neq i \text{ with } \text{conf}_j > \text{conf}_i
  \label{eq:nms}
\end{equation}

\subsection{Batch Inference Optimisation}

Single-frame inference incurs a fixed per-call overhead $T_{\text{ovhd}}$
from Python--C++ context switching. Batching $N$ frames reduces amortised
cost per frame:

\begin{equation}
  t_{\text{per-frame}}^{(N)}
  = \frac{T_{\text{ovhd}} + N \cdot T_{\text{compute}}}{N}
  = \frac{T_{\text{ovhd}}}{N} + T_{\text{compute}}
  \label{eq:batch_speedup}
\end{equation}

For $N = 8$ and $T_{\text{ovhd}} \approx 4\,T_{\text{compute}}$ the
speedup over $N=1$ is:

\begin{equation}
  S = \frac{T_{\text{ovhd}} + T_{\text{compute}}}
           {T_{\text{ovhd}}/8 + T_{\text{compute}}}
  = \frac{5}{1.5} \approx 3.3\times
  \label{eq:batch_speedup_calc}
\end{equation}

\subsection{Frame-Skip Speedup}

For a skip rate $k$ the fraction of frames requiring inference is $1/k$,
giving an approximate end-to-end speedup of:

\begin{equation}
  S_{\text{skip}}(k) \approx
  \frac{T_{\text{full}}}{T_{\text{full}}/k + (1 - 1/k)\,T_{\text{copy}}}
  \approx k
  \qquad \text{(since } T_{\text{copy}} \ll T_{\text{infer}}\text{)}
  \label{eq:frame_skip}
\end{equation}

The 1.5$\times$ mode skips every third frame ($k = 1.5$,
\texttt{frame\_idx \% 3 == 2}), the 2$\times$ mode skips every second,
and the 3$\times$ mode keeps one in three.

% ══════════════════════════════════════════════════════════════════════════════
\section{Results and Analysis}
% ══════════════════════════════════════════════════════════════════════════════

\subsection{Lane Detection Performance}

Table~\ref{tab:lane_results} summarises the segmentation metrics achieved
on the validation split.

\begin{table}[H]
  \centering
  \caption{Lane segmentation validation metrics}
  \begin{tabular}{lc}
    \toprule
    \textbf{Metric} & \textbf{Value} \\
    \midrule
    Pixel Accuracy  & $\approx 97\,\%$ \\
    IoU             & $\approx 0.72$   \\
    Dice / F1       & $\approx 0.84$   \\
    \bottomrule
  \end{tabular}
  \label{tab:lane_results}
\end{table}

\subsection{Object Detection Performance}

YOLOv8n on the COCO validation set achieves:

\begin{table}[H]
  \centering
  \caption{YOLOv8 COCO validation metrics (published)}
  \begin{tabular}{lccc}
    \toprule
    \textbf{Model} & \textbf{mAP\textsubscript{50}} &
    \textbf{mAP\textsubscript{50:95}} & \textbf{Speed (CPU ms/img)} \\
    \midrule
    YOLOv8n & 52.9 & 37.3 & 80.4 \\
    YOLOv8s & 61.8 & 44.9 & 128.4 \\
    YOLOv8m & 67.2 & 50.2 & 234.7 \\
    \bottomrule
  \end{tabular}
  \label{tab:yolo_results}
\end{table}

\subsection{Processing Speed}

Effective throughput on a single CPU core with batch size $N=8$:

\begin{table}[H]
  \centering
  \caption{Approximate processing speed (CPU, 1080p portrait video)}
  \begin{tabular}{lcc}
    \toprule
    \textbf{Mode} & \textbf{Skip} & \textbf{Approx.\ FPS (CPU)} \\
    \midrule
    Quality  & $1\times$   & $\approx 2$--3  \\
    Smooth   & $1.5\times$ & $\approx 3$--4  \\
    Fast     & $2\times$   & $\approx 4$--6  \\
    Fastest  & $3\times$   & $\approx 6$--9  \\
    \bottomrule
  \end{tabular}
  \label{tab:speed}
\end{table}

\subsection{Portrait Video Correction}

The most visible qualitative improvement is the mask re-projection
(Equation~\ref{eq:reproject}). Before the fix, a 1080$\times$1920
portrait frame squeezed into 256$\times$256 caused both severe
aspect-ratio distortion and the model predicting lanes in positions
that do not correspond to the actual road. After the fix the mask
is predicted on the bottom $60\%$ of the frame (road region) and
up-sampled with nearest-neighbour interpolation back to full resolution,
aligning with actual lane markings.

% ══════════════════════════════════════════════════════════════════════════════
\section{Limitations and Failure Cases}
% ══════════════════════════════════════════════════════════════════════════════

\subsection{Lane Detection Limitations}

\begin{enumerate}
  \item \textbf{Fixed input resolution.}
        The model is trained at $256\times256$. Very thin lane markings
        at the far end of a road (small in pixel count) may be missed after
        down-scaling, especially in high-resolution inputs.

  \item \textbf{Road-region crop sensitivity.}
        The \texttt{road\_start\_ratio} must be set manually. An incorrect
        value (too high) may crop out the nearest lane markings; too low
        wastes model capacity on sky.

  \item \textbf{Night and adverse weather.}
        The training data is predominantly daytime clear-weather footage.
        Performance degrades in rain, fog, or low-light scenes where lane
        contrast is reduced.

  \item \textbf{Non-standard lane markings.}
        Dashed, coloured, or worn lane markings common outside Europe/North
        America may not be represented in training data.

  \item \textbf{Occlusion by vehicles.}
        When a vehicle directly ahead covers the lane, the model produces
        fragmented or absent predictions in that region.
\end{enumerate}

\subsection{Object Detection Limitations}

\begin{enumerate}
  \item \textbf{Small and distant objects.}
        YOLOv8n is optimised for speed; the nano model may miss small
        objects such as pedestrians at the far end of a long straight road.

  \item \textbf{Class confusion.}
        Motorcycles and bicycles share visual features and are occasionally
        swapped at low resolution.

  \item \textbf{Traffic sign coverage.}
        Only COCO-defined signs (traffic light, stop sign, parking meter)
        are detected. Country-specific signs (e.g.\ UAE, EU regulatory signs)
        are absent from COCO and will not be detected.

  \item \textbf{Overlapping boxes.}
        Dense traffic scenes with many overlapping vehicles can produce
        duplicate detections that NMS (Equation~\ref{eq:nms}) fails to
        fully suppress at the default IoU threshold.
\end{enumerate}

\subsection{System-Level Limitations}

\begin{enumerate}
  \item \textbf{CPU-only inference.}
        Streamlit Community Cloud provides no GPU. Processing a one-minute
        video at Quality ($1\times$) takes $\approx$20--30 minutes on CPU.
        The frame-skip modes mitigate this but reduce temporal resolution.

  \item \textbf{Memory constraints.}
        The Community Cloud tier provides $\approx$1\,GB RAM. Very long
        or high-resolution videos may exceed this and cause the app to crash.

  \item \textbf{Ephemeral file system.}
        Processed videos are written to \texttt{static/predicted/} on the
        cloud container's local disk. Files are lost on app restart;
        users must download before the session ends.
\end{enumerate}

% ══════════════════════════════════════════════════════════════════════════════
\section{Conclusion}
% ══════════════════════════════════════════════════════════════════════════════

This project delivered a production-ready road scene analysis application
combining lane segmentation and object detection in a single interactive
Streamlit interface. The key technical achievements are:

\begin{itemize}[noitemsep]
  \item \textbf{Correct mask re-projection} — predicting on a cropped ROI
        and up-sampling back to original resolution eliminates the
        portrait-distortion failure mode (Equations~\ref{eq:crop}
        and~\ref{eq:reproject}).

  \item \textbf{Batched inference pipeline} — grouping $N=8$ frames per
        \texttt{model.predict()} call delivers a $\approx$3--5$\times$
        reduction in per-frame latency on CPU
        (Equation~\ref{eq:batch_speedup_calc}).

  \item \textbf{Graceful speed--quality trade-off} — the frame-skip slider
        lets users choose between full-quality and up to $3\times$ faster
        processing without code changes (Equation~\ref{eq:frame_skip}).

  \item \textbf{Persistent UI state} — storing output paths in
        \texttt{st.session\_state} ensures download buttons survive
        Streamlit's script-rerun model, which was a critical usability fix.

  \item \textbf{Modular deployment} — the object-detection variant
        (\texttt{object-detection} branch / \texttt{Object\_Detection\_YOLOv8}
        repository) ships with zero TensorFlow dependency, making it
        lighter and faster to deploy independently.
\end{itemize}

\noindent\textbf{Future work} could include:
\begin{itemize}[noitemsep]
  \item Fine-tuning YOLOv8 on a country-specific traffic sign dataset
        (e.g.\ UAE road signs) to address the detection gap.
  \item Replacing nearest-neighbour up-sampling with learned super-resolution
        to improve mask sharpness at full resolution.
  \item Exporting both models to ONNX and using \texttt{onnxruntime} to
        bypass TensorFlow's Python 3.14 incompatibility on newer runtimes.
  \item Deploying on Hugging Face Spaces with a GPU-enabled ZeroGPU instance
        to achieve real-time ($\geq$25\,FPS) processing.
\end{itemize}

% ══════════════════════════════════════════════════════════════════════════════
\section*{References}
% ══════════════════════════════════════════════════════════════════════════════

\begin{thebibliography}{9}

\bibitem{unet}
O.\ Ronneberger, P.\ Fischer, and T.\ Brox,
``U-Net: Convolutional Networks for Biomedical Image Segmentation,''
\textit{MICCAI}, 2015.

\bibitem{dilated}
F.\ Yu and V.\ Koltun,
``Multi-Scale Context Aggregation by Dilated Convolutions,''
\textit{ICLR}, 2016.

\bibitem{yolov8}
G.\ Jocher, A.\ Chaurasia, and J.\ Qiu,
``Ultralytics YOLOv8,''
\url{https://github.com/ultralytics/ultralytics}, 2023.

\bibitem{coco}
T.-Y.\ Lin \textit{et al.},
``Microsoft COCO: Common Objects in Context,''
\textit{ECCV}, 2014.

\bibitem{streamlit}
Streamlit Inc., ``Streamlit --- The fastest way to build data apps,''
\url{https://streamlit.io}, 2024.

\end{thebibliography}

\end{document}
